Wir nehmen an, das Gewässer bestehe aus Wasser ($\rho_\text{W} = \rhoWO$, FoTaBe S.~187). Auf dem Mond ist die Fallbeschleunigung $g_\text{M} = \gMoO$ (FoTaBe S.~210), der Luftdruck auf der Erde ist etwa der Normaldruck $p_\text{L} = \pLO$ (FoTaBe S.~186). Aus $\rho_\text{W}\,g_\text{M}\,h = p_\text{L}$ folgt
\begin{align}
 h &= \hMoF \\
   &= \frac{\pL}{\rhoW\cdot\gMo} \\
   &= \hMo \approx \result{\hMoP{0}{2}}
\end{align}
Auf der Erde ($g = \gEO$) genügen dafür $\hErF = \hErP{0}{2}$ Wasser; auf dem Mond braucht es wegen der kleineren Fallbeschleunigung etwa sechsmal mehr.
